\documentclass[11pt,a4paper]{ctexart}
\usepackage[top=2.5cm,bottom=2cm,left=2.2cm,right=2.2cm]{geometry}
\usepackage{graphicx,xcolor,titlesec,fancyhdr,hyperref,longtable,booktabs,array,enumitem,tikz,lastpage}
\definecolor{mainred}{HTML}{E94560}
\definecolor{graytext}{HTML}{666666}
\hypersetup{colorlinks=true,linkcolor=mainred,urlcolor=mainred}
\titleformat{\section}{\color{mainred}\Large\bfseries}{\thesection}{1em}{}
\titleformat{\subsection}{\color{mainred}\large\bfseries}{\thesubsection}{1em}{}
\setlist{nosep,leftmargin=*}
\pagestyle{fancy}\fancyhf{}
\fancyhead[L]{\color{graytext}\small 中性原子量子计算}
\fancyhead[R]{\color{graytext}\small 战略情报与成果专利室}
\fancyfoot[C]{\color{graytext}\small ---\ \thepage\ ---}
\renewcommand{\headrulewidth}{0.4pt}
\renewcommand{\headrule}{\hbox to\headwidth{\color{mainred}\leaders\hrule height \headrulewidth\hfill}}
\sloppy
\title{\vspace{2cm}{\Huge\bfseries\color{mainred} 中性原子量子计算发展历程\\[6pt] 与QuEra技术脉络梳理}\\[16pt]\rule{0.5\textwidth}{1.5pt}\\[12pt]{\large 里德堡阻塞 $\cdot$ 相干搬运 $\cdot$ Transversal STAR $\cdot$ 容错量子计算}\\[8pt]{\footnotesize\color{graytext}数据截止：2026年6月}\\[8pt]{\large 战略情报与成果专利室}\\[16pt]{\large\color{graytext}\today}}
\date{}
\begin{document}
\maketitle\thispagestyle{empty}
\newpage\tableofcontents\newpage

\part{中性原子量子计算发展历程}

\section{理论奠基期（2000--2009）：里德堡阻塞与量子门方案}

中性原子路线的理论根基在2000年前后即已奠定。里德堡原子（主量子数$n$很大的激发态原子）具有巨大的电偶极矩，相邻原子间产生强烈的偶极-偶极相互作用，使得在特定空间范围内只能激发一个里德堡原子——这一现象被称为里德堡阻塞效应（Rydberg Blockade）。该阶段还处于纯理论工作，尚未有单原子级别的实验验证。

\textbf{关键论文：}
\begin{itemize}
\item \textbf{Jaksch et al., 2000, Phys. Rev. Lett.}：首次提出利用里德堡阻塞实现快速中性原子量子门的理论方案，奠定了该路线的物理基础。
\item \textbf{Lukin et al., 2001, Phys. Rev. Lett.}：提出``偶极阻塞与介观原子系综中的量子信息处理''，将里德堡阻塞从两体推广至多体系综编码，预言了集体编码量子寄存器的可能性。
\item \textbf{Saffman, Walker \& M\o lmer, 2010, Rev. Mod. Phys.}：系统综述了里德堡原子量子信息处理的完整理论框架，成为该领域引用率最高的综述文献。
\end{itemize}

\section{实验验证期（2009--2017）：从单原子操控到两比特纠缠}

2009年是中性原子路线的实验元年。光镊技术（1986年Ashkin发明）与单原子荧光探测技术的成熟，使得在超高真空腔中捕获、成像、操控单个中性原子成为可能。各学术组独立验证物理可行性，系统误差主要来源于里德堡态有限寿命、光镊加热、原子温度等。比特数从2个扩展到$\sim$50个，但量子门保真度尚不足以支撑容错计算。

\textbf{关键里程碑与论文：}
\begin{itemize}
\item \textbf{Urban et al., 2009, Nature Physics}：威斯康星大学Saffman组首次观测到两个原子间的里德堡阻塞效应。
\item \textbf{Ga\"etan et al., 2009, Nature Physics}：法国Institut d'Optique（Browaeys/Grangier组）独立观测两个原子在里德堡阻塞区的集体激发。
\item \textbf{Wilk et al., 2010, Phys. Rev. Lett.}：Browaeys组利用里德堡阻塞实现两个独立中性原子的纠缠，保真度约70\%，标志着中性原子两比特门的首次实验演示。
\item \textbf{Isenhower et al., 2010, Phys. Rev. Lett.}：Saffman组演示了中性原子受控非门（CNOT），进一步验证了量子逻辑操作的可行性。
\item \textbf{Bernien et al., 2017, Nature}：哈佛Lukin组实现51原子可编程量子模拟器，利用光镊阵列制备一维链和二维晶格，观测多体动力学，首次展示中性原子在``量子模拟''而非``量子计算''方向的规模优势。
\end{itemize}

\section{规模扩展期（2017--2021）：从模拟器到可编程阵列}

此阶段的核心成就是将原子阵列规模从数十个推向数百个，并引入原子重排（rearrangement）技术——通过移动光镊将原子重新排列成无缺陷的规则阵列，解决了``光镊加载随机性导致部分格点空缺''的问题。中性原子开始被视为``量子模拟''的有力工具，但量子门操作仍受限于近邻连接——里德堡门只能在空间相邻的原子对之间执行。远距离原子间的纠缠需要通过一系列SWAP操作逐跳传递，效率极低。

\textbf{关键论文：}
\begin{itemize}
\item \textbf{Levine et al., 2019, Phys. Rev. Lett.}：哈佛Lukin组提出并演示并行高保真多量子比特门，通过全局激光脉冲同时驱动多对原子，大幅提升了门操作的并行度。
\item \textbf{Ebadi et al., 2021, Nature}：哈佛/QuEra团队实现256原子可编程量子模拟器，在二维光镊阵列中制备了量子自旋液体态、测度了纠缠熵，展示了中性原子在``规模''上的绝对优势。这是当时全球最大规模的可控量子系统之一。
\end{itemize}

\section{架构突破期（2022）：从``近邻连接''到``任意连接''}

2022年是中性原子路线的关键转折年。4月，哈佛/QuEra团队在Nature发表了一篇具有范式意义的论文，彻底改变了该路线的架构前景。这篇论文发表后，投资界和学术界对中性原子的定位发生质变：从``第三条路线的有趣补充''升级为``可能在纠错效率上超越超导的通用计算平台''。

\textbf{关键论文：}
\begin{itemize}
\item \textbf{Bluvstein et al., 2022, Nature}：``A quantum processor based on coherent transport of entangled atom arrays''。该工作首次实现了在保持量子相干的前提下，用光镊将已经纠缠好的原子团整体搬运到目标位置，再与另一组原子执行门操作。这解决了此前中性原子路线最核心的架构瓶颈——连接性受限。具体演示了：（1）可编程生成团簇态（cluster states）、Steane码（7-qubit）；（2）制备表面码态（13 data + 6 ancilla）和环面码态（16 data + 8 ancilla）；（3）混合模拟-数字演化，测量量子多体疤痕的纠缠熵非单调动力学。
\item \textbf{Graham et al., 2022, Nature}：威斯康星大学Saffman组同期发表``Multi-qubit entanglement and algorithms on a neutral-atom quantum computer''，展示了在中性原子平台上运行多比特纠缠和量子算法的能力，进一步证明了该路线的通用性。
\end{itemize}

\begin{table}[h]
\centering
\caption{2022年架构突破的前后对比}
\begin{tabular}{p{2.5cm}p{5cm}p{5cm}}
\toprule
\textbf{维度} & \textbf{此前的局限} & \textbf{2022年的改变} \\
\midrule
连接性 & 仅限近邻，像超导一样受困于固定布局 & 原子可移动，实现全连接/任意连接 \\
\midrule
纠错码适配 & 表面码需要大量长距离辅助比特操作 & 直接演示了表面码、环面码制备，证明中性原子天然适合LDPC \\
\midrule
架构范式 & 中性原子被视为量子模拟器，缺乏通用计算架构 & 首次明确展示了通用量子处理器架构路径 \\
\midrule
产业展望 & 扩展前景不明，可能止步于学术验证 & 确定性资本涌入，规模与效率同步提升 \\
\bottomrule
\end{tabular}
\end{table}

\section{纠错与逻辑量子比特期（2023--2024）：从物理比特到逻辑比特}

2022年的架构突破为纠错演示铺平了道路。2023--2024年，中性原子路线在纠错领域取得一系列里程碑式进展，迅速缩小与超导路线的差距。技术竞争从``比物理比特数''转向``比纠错架构效率''和``比逻辑比特产出率''。中性原子的动态全连接优势在纠错场景中开始显现。

\textbf{关键论文：}
\begin{itemize}
\item \textbf{Evered et al., 2023, Nature}：哈佛/QuEra团队演示高保真并行纠缠门，单比特门保真度达99.9\%+，两比特门保真度达99.5\%+，逼近超导水平。
\item \textbf{Scholl et al., 2023, arXiv}：提出并演示擦除转换（erasure conversion）技术，将物理错误转换为可定位的擦除错误，大幅提升了纠错效率。
\item \textbf{Bluvstein et al., 2024, Nature}：``Logical quantum processor based on reconfigurable atom arrays''。这是中性原子路线的里程碑论文——首次基于可重构原子阵列构建了逻辑量子处理器，演示了：（1）多达48个逻辑量子比特的纠缠；（2）利用原子移动实现跨逻辑量子比特的横向操作（transversal operations）；（3）逻辑级别的纠错与计算。
\item \textbf{Reichardt et al., 2024, arXiv}：在中性原子量子处理器上演示逻辑计算，进一步验证了该路线在容错计算上的可行性。
\end{itemize}

\section{架构创新与产业化加速期（2025--2026）：STAR架构与Megaquop时代}

2025--2026年，中性原子路线进入``架构级创新''的爆发期，核心目标是降低容错量子计算的物理资源开销，向``megaquop''（百万级可靠逻辑操作）乃至``teraquop''时代迈进。中性原子路线开始展现``算法-纠错码-硬件''协同设计的独特优势。其动态全连接能力不仅是一种工程便利，更成为降低容错开销的核心物理资源。

\textbf{关键论文：}
\begin{itemize}
\item \textbf{Xu et al., 2024, Nature Physics}：提出在可重构原子阵列上实现qLDPC码的硬件高效方案。研究表明，中性原子平台仅需数百个物理比特即可开始超越表面码，将物理-逻辑比特比从$\sim$1000:1压缩到低十位。
\item \textbf{Ismail et al., 2026, PRX Quantum}：QuEra与洛斯阿拉莫斯国家实验室联合发表``Transversal architecture for megaquop-scale quantum simulation with neutral atoms''，提出Transversal STAR架构。该架构：（1）通过横向注入（transversal injection）直接制备小角度magic state，消除了传统方案中的Solovay-Kitaev合成开销；（2）利用中性原子的可重构连接性和大规模并行性，将Clifford门操作时间匹配模拟旋转速度；（3）在表面码下，仅需约10,000物理比特即可实现模拟体积$>600$的哈密顿演化；（4）结合高码率qLDPC码，物理比特数可进一步降至1,500--3,000；（5）相比固定连接架构，执行速度提升约250倍。
\item \textbf{Zhao et al., 2026, arXiv}：演示码率超过1/2的超高码率qLDPC码，在0.1\%物理错误率下实现每逻辑比特每轮约$1.3\times10^{-13}$的逻辑错误率，逼近``teraquop''（万亿级可靠操作） regime。
\end{itemize}

\part{QuEra Computing}

\section{公司概况}

QuEra Computing是一家专注于中性原子量子计算的美国量子科技公司，2018年由哈佛大学与麻省理工学院（MIT）的顶尖量子物理团队联合创立，总部位于美国马萨诸塞州波士顿。公司从中性原子量子计算路线的奠基性学术成果中孵化而来，是目前全球该领域融资规模最大、技术路线最清晰、产业化进程最快的企业。

QuEra已构建``硬件+软件+云服务''的全栈产品体系：
\begin{itemize}
\item \textbf{Aquila}：256量子比特中性原子量子计算机（2022年商用），支持模拟和数字双模式运算，是AWS Braket上唯一的中性原子量子系统
\item \textbf{Bloqade}：开源量子编程SDK，支持Python接口和脉冲级控制，专门为中性原子的模拟和数字混合模式优化
\item \textbf{Tsim}：快速纠错模拟器（2026年发布），支持非Clifford电路的通用模拟
\item \textbf{Gemini-Class}：下一代硬件，目标支持横向逻辑门和qLDPC编码，部署于日本AIST的ABCI-Q超算中心
\end{itemize}

\subsection{融资历程}

QuEra是中性原子路线中融资能力最强的企业，累计融资约\$3.3亿美元，投资方涵盖科技巨头、主权资本与顶级风投：
\begin{itemize}
\item 2021年11月：\$1,700万A轮（Rakuten、Day One Ventures等）
\item 2022年8月：DARPA ONISQ计划合同\$1,100万
\item 2023年6月：\$3,600万B轮（Eclipse Ventures领投）
\item 2024年11月：日本AIST \$4,100万国家级部署合同
\item 2025年2月：\textbf{\$2.3亿战略轮}（Google Quantum AI与软银愿景基金联合领投），估值推至约\$15亿
\end{itemize}

\$2.3亿的规模在中性原子路线中遥遥领先——同路线对手Pasqal累计约\$1.4亿，Atom Computing已被Quantinuum吸收。更关键的是，这轮融资后QuEra的论文产出从2025年中的6篇密集爆发直接推进到2026年的Libra路线图，资金$\rightarrow$研发$\rightarrow$产品的时间压缩效应极为明显。

\subsection{商业化布局}

\begin{itemize}
\item \textbf{云服务}：Aquila作为AWS Braket上唯一的中性原子量子计算机，已面向全球开发者开放；2026年6月与AWS联合宣布Libra 2028部署计划
\item \textbf{企业合作}：与Deloitte、BCG、NTT DATA、SAS等建立行业应用联盟，覆盖医疗、金融、优化等领域
\item \textbf{HPC融合}：与Pawsey超算中心、NERSC、HPE等推进量子-经典混合计算架构
\item \textbf{全球化}：波士顿总部 + 日本AIST部署 + 英国Harwell测试床 + 澳洲Pawsey合作
\end{itemize}

\section{创始团队}

QuEra的创始团队具有``四位顶尖教授+连续创业高管''的独特组合：

\begin{itemize}
\item \textbf{Mikhail Lukin}（联合创始人兼首席科学家）：哈佛大学物理学教授，美国国家科学院院士，h-index 176，450+论文，136,000+引用。中性原子量子计算路线的奠基人，Transversal STAR架构和qLDPC纠错方案的核心思想均源自其哈佛实验室。

\item \textbf{Vladan Vuleti\'c}（联合创始人兼CTO）：MIT Lester Wolfe物理学教授，h-index 65，236篇论文。2017年发表标志性的``51原子量子模拟器''论文（Nature, 3,190引用），为QuEra硬件平台奠定直接基础。师从诺贝尔奖得主Steven Chu。

\item \textbf{Markus Greiner}（联合创始人）：哈佛大学物理学教授，h-index 55，麦克阿瑟天才奖获得者。开发了光晶格中单原子分辨成像技术——中性原子量子计算机可编程性的前提技术。

\item \textbf{Dirk Englund}（联合创始人）：MIT EECS教授，h-index 81，725篇论文。同为Lightmatter（光子AI芯片，估值数十亿）、DUST Identity等的联合创始人，是将量子光子学从实验室推向产业化的跨界典范。

\item \textbf{Nathan Gemelke}（联合创始人兼首席技术战略官）：斯坦福大学物理学博士（导师为诺贝尔奖得主Steven Chu），宾州州立大学前助理教授，Alfred P. Sloan研究学者。

\item \textbf{Andy Ory}（CEO）：连续创业者。先后创办Priority Call Management（以\$1.6亿出售）、Acme Packet（IPO 2006，峰值市值\$40亿，被Oracle以\$21亿收购）、128 Technology（被Juniper收购）。

\item \textbf{Takuya Kitagawa}（总裁）：哈佛大学物理学博士，前Rakuten AI部门负责人（从1人扩展到1,000+人团队）。2024年11月加入QuEra，仅4个月后即主导完成Google/软银\$2.3亿战略融资。
\end{itemize}

\section{技术脉络}

\subsection{第一阶段：学术孵化与创业奠基（2018--2021）}

QuEra的技术起点深植于哈佛大学Lukin组的学术积累，但公司成立初期的核心任务是将实验室的``量子模拟器''转化为可工程化的计算平台。

QuEra于2018年由Mikhail Lukin、Vladan Vuleti\'c、Markus Greiner、Dirk Englund及Nathan Gemelke联合创立。创始团队全部来自哈佛/MIT的顶尖AMO实验室，其技术遗产包括2016年的逐原子组装技术、2017年的51原子量子模拟器，以及2019年由Harry Levine提出的并行高保真多量子比特门。这些成果为QuEra提供了两项关键资产：光镊阵列的规模化制备能力和高保真并行门操控技术。

然而，QuEra成立之初面临一个结构性困境：中性原子平台虽然比特数可扩展（2021年Ebadi等人在哈佛实现了256原子阵列），但量子门操作受限于近邻连接，远距离原子间的纠缠需要逐跳SWAP传递，效率极低。这意味着QuEra继承的硬件平台本质上仍是``量子模拟器''——适合制备特定量子态并观测多体动力学，但缺乏通用量子计算所需的任意连接能力。此时，公司的商业化路径尚不清晰，Alex Keesling（创始CEO）带领团队完成了A轮融资（\$1,700万，2021年），但产品形态仍停留在学术合作与原型验证阶段。

这一阶段QuEra的核心矛盾是：拥有规模化物理平台，但缺乏定义其``计算架构''的关键突破。

\subsection{第二阶段：架构定义与Aquila商用化（2022--2023）}

2022年，QuEra通过一项关键学术合作获得了定义自身技术架构的``范式级突破''，并在此后一年内完成了从``原型机''到``商用产品''的跨越。

2022年4月，哈佛Lukin组博士研究生Dolev Bluvstein（QuEra核心合作者，后获Physics World 2024年度突破奖）与Harry Levine（QuEra关联研究者）等人在Nature发表了``A quantum processor based on coherent transport of entangled atom arrays''。该工作首次实现了相干搬运——在保持量子纠缠的前提下，用光镊将原子团整体移动到目标位置。这一技术彻底改变了QuEra的架构前景：光镊从``固定原子的工具''转变为``动态布线''的手段，使得任意远距离量子比特间的逻辑门成为可能，而无需增加物理连接。

对QuEra而言，这一突破具有决定性意义。此前，中性原子被视为超导和离子阱之外的``第三条路线''，一个连接受限的``量子模拟器''补充；相干搬运之后，QuEra可以宣称其平台同时具备``可扩展性''和``全连接性''——这是当时任何其他技术路线都无法提供的组合。公司迅速将这一学术成果转化为产品基础：2022年11月，QuEra发布了Aquila——256量子比特中性原子量子计算机，成为全球首款上架AWS Braket的中性原子量子系统。

2023年，QuEra进一步巩固了其技术领先性。Simon Evered（QuEra研究科学家）等人在Nature上展示了高保真并行纠缠门，单比特保真度99.9\%+、两比特保真度99.5\%+，基本追平超导水平，为容错计算提供了物理基础。紧接着，Bluvstein与Evered、Hengyun Zhou（QuEra量子纠错架构负责人，后任MIT助理教授）等人在同年的Nature上发表了基于可重构原子阵列的逻辑量子处理器——在280个物理量子比特上实现48个逻辑量子比特的纠缠，并展示了跨逻辑比特的横向操作。这是QuEra首次在``逻辑层面''证明其平台的计算能力，标志着公司从``硬件供应商''向``容错计算平台''转型。

这一阶段QuEra的战略成果是：以相干搬运为架构核心，以Aquila为产品载体，以逻辑处理器为技术里程碑，确立了中性原子路线的产业化标杆。

\subsection{第三阶段：纠错深化与路线分化（2023--2024）}

在证明逻辑比特可行性后，QuEra面临一个更深层的技术抉择：选择哪种纠错码架构？这一抉择最终将公司推向``双轨并行''战略。

2024年，Qian Xu（QuEra关联研究者）等人在Nature Physics上证明了可重构原子阵列可以实现``恒定开销''容错量子计算——即物理-逻辑比特比不随系统规模膨胀。这一结果对QuEra极具吸引力，因为它意味着中性原子的动态全连接能力可以转化为纠错效率优势，而非仅仅是工程便利。然而，QuEra团队在深入探索qLDPC码（量子低密度奇偶校验码）时发现了一个结构性限制：作为CSS型码，qLDPC可以横向实现CNOT门，但通常无法横向实现S门（相位门）。这意味着连完整的Clifford群都无法在编码块内具备，更遑论通用量子计算所需的非Clifford门（如T门）。Cain（QuEra研究科学家）等人在2024年提出的横向门关联解码策略，只能在现有框架下优化效率，无法从根本上解决门集不完整的问题。

这一发现迫使QuEra重新审视其技术路线图。公司意识到，若坚持在qLDPC码上追求通用容错计算，将不得不依赖Magic State蒸馏等非横向手段，这会大幅削弱中性原子动态全连接带来的``低开销''优势；反之，若接受qLDPC的横向门限制，则可以利用其高码率特性，在特定结构的量子模拟任务上实现数量级的资源节省。2024年底至2025年初，QuEra做出了关键的战略分化：不再将单一赌注押在``通用FTQC''上，而是同时开辟一条``专用量子模拟''的近期路线。

这一阶段QuEra的核心决策是：承认qLDPC横向门集的结构性限制，将技术路线从``单轨通用''调整为``双轨并行''——一条轨道继续向通用容错计算推进，另一条轨道利用限制本身，在专用模拟领域建立近期优势。

\subsection{第四阶段：双轨战略——STAR架构与通用FTQC（2025--2026）}

2025年至2026年，QuEra的双轨战略全面展开：以Transversal STAR架构主攻专用量子模拟，以96逻辑比特演示和Libra系统路线图推进通用容错计算。

\textbf{轨道一：专用模拟——STAR架构。} 2025年，Refaat Ismail（QuEra研究科学家）与Zhou（前QuEra QEC架构负责人，MIT助理教授）、Zhao（QuEra研究科学家）、Liu（QuEra高级研究科学家）、Wang（QuEra VP，算法与应用）等人在arXiv上提出、并于2026年在PRX Quantum正式发表了Transversal STAR架构。该架构明确接受其qLDPC码横向门集不包含S门的事实，将自身定位为``结构化量子模拟专用架构''。它利用中性原子的动态全连接，通过``横向注入''直接制备小角度旋转态，配合后选择实现部分容错的模拟旋转门，从而绕过了Magic State蒸馏和Solovay-Kitaev合成的高开销。在表面码保护下，STAR仅需约一万个物理量子比特即可实现模拟体积超过600的哈密顿演化；若结合高码率qLDPC码，物理比特数可降至1,500--3,000个。QuEra官方明确承认STAR的专用性：它不是通用量子计算机，而是``通往megaquop级量子模拟的最便宜路径''。

\textbf{轨道二：通用容错——Libra路线图。} 与此同时，QuEra并未放弃通用FTQC。Zhou（前QuEra）与Zhao（QuEra）等人在2025年Nature上发表的``低开销横向容错''方案，以及同年11月Bluvstein（QuEra核心合作者）与Zhou（前QuEra）、Zhao（QuEra）等人展示的96逻辑量子比特演示（低于阈值的$d=5$表面码纠错，通用门集$\{H, T, CNOT\}$），代表了中性原子通用容错路线的最高水平。基于此，QuEra宣布了与AWS的战略合作：2028年在Amazon Braket上线Libra系统，目标提供256逻辑量子比特的纠错级量子计算云服务。

在纠错码研究层面，Zhao（QuEra）等人在2026年的arXiv预印本中展示了码率超过1/2的超高码率qLDPC码，逼近``teraquop''阈值；同年Haenel、Luo与Zhao（QuEra）发布的Tsim开源纠错模拟器，则为双轨战略提供了统一的验证工具链。

这一阶段QuEra的战略格局是：以STAR架构在近期量子模拟领域证明量子优势、获取商业营收；以通用FTQC路线（Libra系统）在中长期建立技术壁垒。两条轨道共享同一套中性原子硬件基础（光镊阵列、相干搬运、高保真门），但在纠错码架构和门实现方式上分道扬镳。

\section{关键论文汇总}

\begin{longtable}{p{2cm}p{1.2cm}p{3cm}p{5cm}p{4cm}}
\toprule
\textbf{阶段} & \textbf{年份} & \textbf{作者} & \textbf{期刊} & \textbf{论文贡献} \\
\midrule
\endhead
\midrule
\multicolumn{5}{r}{\footnotesize 续下页}
\endfoot
\bottomrule
\endlastfoot
理论奠基 & 2000 & Jaksch et al. & PRL 85, 2208 & 首次提出基于里德堡阻塞的量子门方案 \\
理论奠基 & 2001 & Lukin et al. & PRL 87, 037901 & 系综中的偶极阻塞与量子信息处理 \\
理论奠基 & 2010 & Saffman et al. & RMP 82, 2313 & 领域引用率最高的综述文献 \\
\midrule
实验验证 & 2009 & Urban et al. & Nat. Phys. 5, 110 & 首次观测两个原子间里德堡阻塞 \\
实验验证 & 2010 & Wilk et al. & PRL 104, 010502 & 首次演示两个中性原子的纠缠（$\sim$70\%） \\
实验验证 & 2010 & Isenhower et al. & PRL 104, 010503 & 演示中性原子CNOT门 \\
实验验证 & 2017 & Bernien et al. & Nature 551, 579 & 51原子可编程量子模拟器 \\
\midrule
规模扩展 & 2019 & Levine et al. & PRL 123, 170503 & 并行高保真多比特门 \\
规模扩展 & 2021 & Ebadi et al. & Nature 595, 227 & 256原子可编程模拟器，全球最大 \\
\midrule
架构突破 & 2022 & Bluvstein et al. & Nature 604, 451 & 相干搬运：原子阵列的范式级突破 \\
架构突破 & 2022 & Graham et al. & Nature 604, 457 & 中性原子多比特纠缠与算法 \\
\midrule
纠错与逻辑 & 2023 & Evered et al. & Nature 622, 268 & 高保真并行纠缠门（99.9\%+） \\
纠错与逻辑 & 2024 & Bluvstein et al. & Nature 626, 58 & \textbf{逻辑量子处理器：48逻辑比特} \\
\midrule
STAR架构 & 2024 & Xu et al. & Nat. Phys. 20, 832 & 恒定开销容错量子计算 \\
STAR架构 & 2025 & Bluvstein et al. & Nature (2025) & 96逻辑比特，通用门集演示 \\
STAR架构 & 2026 & Ismail et al. & PRX Quantum & \textbf{Transversal STAR架构（Los Alamos合作）} \\
STAR架构 & 2026 & Zhao et al. & arXiv & 超高码率qLDPC（$10^{-13}$错误率） \\
\end{longtable}

\section{展望与评估}

QuEra在2029年实现1000逻辑量子比特的目标技术上可行但不确定性显著：其优势在于已硬件验证48$\rightarrow$96逻辑比特的below-threshold纠错、3000原子连续运行及魔法态蒸馏等核心模块，且2025年Nature论文已展示通用门集$\{H, T, CNOT\}$的96逻辑比特演示。但关键瓶颈在于通用路线对非Clifford门（T门）的依赖——qLDPC码无法横向实现S门/T门，必须通过魔法态蒸馏补偿，而蒸馏开销随系统规模超线性增长。同时qLDPC码仅停留在模拟验证阶段，若硬件验证失败将被迫回归表面码，物理/逻辑比从20:1恶化至100:1，使1000通用逻辑比特需要10万+物理量子比特。2027年Libra原型机的通用门集验证与qLDPC硬件验证将成为决定性分水岭。

\end{document}
